\chapter{Introduction}
\label{chapter:introduction}

Personal agents are becoming delegates: systems that hold private data, wield tools, and act on behalf of their owners. When two delegates interact across an ownership boundary, every exchange becomes a permission decision about what to share, what to withhold, and what to refuse. This thesis investigates the security and utility consequences of cross-boundary agent delegation, proposes benchmarks to measure them, and develops architectural solutions that shift enforcement from probabilistic reasoning to deterministic structure.

The central observation is simple. A delegate that refuses every cross-boundary request is perfectly safe but entirely useless. A delegate that answers every request is maximally useful but catastrophically unsafe. Between these extremes lies a frontier of achievable operating points. We call this the \textit{security-utility frontier}. The shape of this frontier, the factors that move it, and the architectural interventions that can expand it are the subject of this thesis.


% ============================================================================
\section{The Rise of Personal Agent Delegates}
\label{sec:intro_rise}
% ============================================================================

The trajectory from language model to personal delegate has unfolded in three phases. \textbf{Phase One (Tools)} transformed language models into tool-users: ReAct~\cite{yao2023react} interleaves reasoning with executable actions, and Toolformer~\cite{schick2023toolformer} showed tool use could be a learned competency. These systems were stateless and had limited attack surface.

\textbf{Phase Two (Assistants)} produced stateful systems. Reflexion~\cite{shinn2023reflexion} enabled agents that critique and improve their own outputs across turns. MemGPT~\cite{packer2024memgpt} introduced an OS-inspired memory hierarchy giving agents persistent long-term storage. An assistant with persistent memory holds a growing corpus of personal information, making it increasingly useful but also increasingly dangerous if exposed to unauthorised parties.

\textbf{Phase Three (Delegates)}, now underway, produces persistent, autonomous systems that act on behalf of a specific person. Devin~\cite{devin2024} operates as an autonomous software engineer. Manus~\cite{manus2025} pushes desktop automation with persistent project state. These are not chatbots. They are digital proxies holding private context and wielding powerful tools.

Two protocol-level developments accelerate the transition to a networked ecosystem. The Model Context Protocol (MCP)~\cite{anthropic2024mcp} standardises how agents connect to tools and data sources, analogous to USB for software. The Agent-to-Agent protocol (A2A)~\cite{google2025a2a} addresses discovery, authentication, and communication between agents. Together, MCP provides the tool layer and A2A the networking layer. But neither addresses the operational question at the heart of every cross-boundary interaction: \textit{what should one agent be permitted to disclose about its owner to another agent?} The infrastructure for connection exists. The infrastructure for safe connection does not.


% ============================================================================
\section{The Cross-Boundary Delegation Problem}
\label{sec:intro_cross_boundary}
% ============================================================================

Consider a system with two agents, Agent~A and Agent~B, acting on behalf of Owner~A and Owner~B respectively. Agent~A holds a context~$C_A$ consisting of Owner~A's private information. Agent~B initiates a request~$r$ that requires some subset of~$C_A$ to be fulfilled. The cross-boundary delegation problem is: Agent~A must decide, for each request~$r$, what portion of~$C_A$ to disclose, balancing \textbf{utility} (sufficient information to fulfil the request) against \textbf{security} (no disclosure Owner~A would not sanction).

As a concrete example: Alice's agent holds her calendar, including an oncology follow-up (Tuesday), an interview at a competing firm (Wednesday), and a therapy session (Thursday). Bob's agent asks to schedule a meeting. The ideal response shares time-slot availability without revealing \textit{why} slots are blocked: ``Tuesday afternoon doesn't work'' rather than ``I can't do Tuesday because of a medical appointment.'' This requires distinguishing \textit{what} information is needed from the sensitive content underneath, a distinction absent from traditional access control policies.

The problem is pervasive~\cite{talebirad2023multiagent, guo2024largelanguagemodelbased}: it arises whenever a colleague's agent asks for a project update (share progress, not internal conflicts), a recruiter's agent probes career interests (openness vs.\ discretion), or a vendor's agent negotiates budget (reveal constraints selectively). Critically, the problem exhibits an asymmetry of error: over-disclosure is irreversible, while under-disclosure can be retried or escalated.


% ============================================================================
\section{Why This Is Different From Traditional Access Control}
\label{sec:intro_different}
% ============================================================================

Traditional access control~\cite{bell1973secure, denning1976lattice} operates on a principle where \textbf{policy is enforcement}: a security label mechanistically determines access. The reference monitor is deterministic and formally verifiable.

LLM-based delegation violates this principle. The privacy policy is not enforcement but \textit{one input to a reasoning process that produces enforcement as an output}. The agent receives the policy, the query, the conversational context, and the relationship, then reasons about what to do. The outcome is stochastic: the same agent, same policy, same query may produce different responses depending on phrasing, preceding conversation, and random seed.

This has three consequences. First, the frontier \textit{cannot be predicted from policy text alone}. Two policies with similar intent may produce different boundaries because the model interprets them differently. Second, the frontier \textit{can only be measured empirically}, which is why we need benchmarks. Third, the frontier \textit{is unstable}: implicit social norms from pretraining compose with explicit policy in ways the policy author cannot anticipate.

These properties mean that formal verification and lattice models are insufficient. We need empirical measurement of emergent boundaries, benchmarks that capture both dimensions simultaneously, and architectural interventions that recover the determinism of traditional access control.


% ============================================================================
\section{Research Questions}
\label{sec:intro_rqs}
% ============================================================================

\begin{tcolorbox}[colback=blue!3!white, colframe=blue!40!black, title={\textbf{RQ1: How do we measure cross-boundary delegation?}}]
How should a personal-agent system represent the boundary between legitimate coordination and privacy leakage? This question motivates the SharedOS problem formulation and the PACT-PAIR dyadic benchmark: before proposing defences, we need an evaluation that measures security and utility together rather than treating refusal as success.
\end{tcolorbox}

\begin{tcolorbox}[colback=blue!3!white, colframe=blue!40!black, title={\textbf{RQ2: What fails when delegation becomes networked?}}]
Which failures appear only when agents interact across a social graph rather than a single pair? This question motivates PACT-NET, where transitive leakage, confused deputies, cross-cluster disclosure, and write-side laundering become measurable network-scale phenomena.
\end{tcolorbox}

\begin{tcolorbox}[colback=green!3!white, colframe=green!40!black, title={\textbf{RQ3: Can architecture move the frontier beyond prompting?}}]
Can enforcement be shifted from model instruction-following into the execution environment itself? This question motivates Mountable Context Cells and the Escalation Protocol: mechanisms that remove unauthorised data before retrieval or stop ambiguous tool calls before private context enters the model.
\end{tcolorbox}


% ============================================================================
\section{Contributions}
\label{sec:intro_contributions}
% ============================================================================

\begin{enumerate}[leftmargin=2em]

    \item \textbf{Problem formulation.} We formalise cross-boundary delegation as emergent enforcement, define the security-utility frontier, and show why traditional access control is insufficient (\Cref{chapter:problem_formulation}).

    \item \textbf{Platform.} We design SharedOS, a multi-agent delegation system with per-relationship context scoping, configurable policies, and tool orchestration. It serves as both a deployed system and experimental apparatus (\Cref{chapter:problem_formulation}).

    \item \textbf{Benchmarks.} PACT-PAIR (600 dyadic scenarios, five sensitivity categories, six-level defence gradient) and PACT-NET (997 network scenarios testing transitive disclosure). Both use dual-metric evaluation measuring utility and security simultaneously (\Cref{chapter:benchmark_design,chapter:failure_cases}).

    \item \textbf{Empirical insights.} We show that generic privacy instructions do not move the frontier, category-specific policies have a ceiling, multi-turn interaction exposes incidental leakage invisible to single-turn tests, and relationship context mainly shifts the cost toward over-refusal (\Cref{chapter:benchmark_design,chapter:failure_cases}).

    \item \textbf{Architectural solutions.} \textit{Relationship-specific policy} as a behavioural layer, \textit{Mountable Context Cells} for structural data absence, and the \textit{Escalation Protocol} to stop ambiguous tool calls before retrieval (\Cref{chapter:solution_proposal}).

\end{enumerate}


% ============================================================================
\section{Thesis Organisation}
\label{sec:intro_organisation}
% ============================================================================

\Cref{chapter:related_work} reviews agent architectures, LLM security, multi-agent protocols, and existing benchmarks. \Cref{chapter:problem_formulation} formally defines cross-boundary delegation and presents the SharedOS platform. \Cref{chapter:benchmark_design} describes the PACT-PAIR benchmark design and reports dyadic experimental results across the defence gradient. \Cref{chapter:failure_cases} presents PACT-NET, extending evaluation to multi-agent network scenarios with transitive disclosure risks. \Cref{chapter:solution_proposal} proposes Mountable Context Cells and the Escalation Protocol, with validation experiments. \Cref{chapter:conclusion_discussion} discusses findings, identifies limitations, and presents directions for future work.
